% Problem-section file following ../_template.tex.
\section{Absolute separability from spectra}
\label{sec:problem-15}

\paragraph{Problem.}
Characterize the spectra of bipartite states that remain separable under every
global unitary, and decide whether absolute separability equals absolute
positivity under partial transpose in higher local dimensions.  Let
$2\le m\le n$ and let
$\lambda=(\lambda_1,\ldots,\lambda_{mn})$ be a decreasing probability vector.
Define the absolutely separable spectra by
\begin{equation}
  \operatorname{ASEP}_{m,n}
  :=\left\{
    \lambda:
    U\operatorname{diag}(\lambda)U^\dagger
    \text{ is separable on }\mathbb{C}^m\otimes\mathbb{C}^n
    \text{ for every }U\in\mathrm{U}(mn)
  \right\}.
  \label{eq:p15-asep-spectra}
\end{equation}
The relaxation by positivity under partial transpose is
\begin{equation}
  \operatorname{APPT}_{m,n}
  :=\left\{
    \lambda:
    \left(U\operatorname{diag}(\lambda)U^\dagger\right)^{T_B}\succeq0
    \text{ for every }U\in\mathrm{U}(mn)
  \right\}.
  \label{eq:p15-appt-spectra}
\end{equation}
The task is to characterize the set in Eq.~\eqref{eq:p15-asep-spectra} and,
for $3\le m\le n$, decide whether it equals the set in
Eq.~\eqref{eq:p15-appt-spectra}.

\subsection*{Status}

Unsolved

\subsection*{Source}

Kr\"uger and Werner pose the spectral characterization of absolute
separability, and Ahiable, Kothakonda, and Winter explicitly retain the
higher-dimensional characterization and ASEP--APPT equality as open
\sourcecite{ref:p15-krueger-werner}{KW05},
\sourcecite{ref:p15-ahiable-et-al}{AKW26}.

\subsection*{Progress}

\begin{itemize}
  \item For two qubits, Verstraete, Audenaert, and De Moor obtained the exact
  condition
  \begin{equation}
    \lambda\in\operatorname{ASEP}_{2,2}
    \quad\Longleftrightarrow\quad
    \lambda_1-\lambda_3-2\sqrt{\lambda_2\lambda_4}\le0.
    \label{eq:p15-two-qubit-criterion}
  \end{equation}
  Equation~\eqref{eq:p15-two-qubit-criterion} settles only the smallest
  bipartite dimension
  \sourcecite{ref:p15-verstraete-et-al}{VAD01}.

  \item Hildebrand gave necessary-and-sufficient linear-matrix-inequality
  conditions for membership in Eq.~\eqref{eq:p15-appt-spectra} in arbitrary
  finite dimensions.  This characterizes the relaxation, not absolute
  separability
  \sourcecite{ref:p15-hildebrand}{Hil07}.

  \item Johnston proved
  $\operatorname{ASEP}_{2,n}=\operatorname{APPT}_{2,n}$ for every $n$, so all
  cases with a qubit factor are excluded from the remaining problem
  \sourcecite{ref:p15-johnston}{Joh13}.

  \item Abellanet-Vidal et al. derived sufficient spectral criteria for
  absolute separability in arbitrary dimensions by inverting positive maps
  and taking convex hulls of the resulting regions.  These provide
  computable inner approximations to
  $\operatorname{ASEP}_{m,n}$, but are not necessary conditions
  \sourcecite{ref:p15-abellanet-vidal-et-al}{AMR+25}.

  \item Ahiable, Kothakonda, and Winter established new geometric properties
  of both spectral sets.  In particular, $\operatorname{APPT}_{m,n}$ is a
  spectrahedron whose faces are all exposed, whereas
  $\operatorname{ASEP}_{m,n}$ is semialgebraic.  They explicitly retain the
  higher-dimensional characterization and equality as open
  \sourcecite{ref:p15-ahiable-et-al}{AKW26}.

  \item Tran subsequently derived an explicit dimension-dependent upper
  bound on the purity of every absolutely-PPT state.  This is a quantitative
  outer constraint on $\operatorname{APPT}_{m,n}$, not a spectral
  characterization, and hence does not settle its equality with
  $\operatorname{ASEP}_{m,n}$
  \sourcecite{ref:p15-tran}{Tra26}.
\end{itemize}

\subsection*{References}

\begin{enumerate}[label={},leftmargin=4.8em,labelsep=0.5em]
  \footnotesize
  \item[\textup{[VAD01]}]\label{ref:p15-verstraete-et-al}
  F. Verstraete, K. Audenaert, and B. De Moor,
  ``Maximally Entangled Mixed States of Two Qubits,''
  \emph{Physical Review A} \textbf{64}, 012316 (2001).
  \href{https://doi.org/10.1103/PhysRevA.64.012316}%
       {doi:10.1103/PhysRevA.64.012316};
  \href{https://arxiv.org/abs/quant-ph/0011110}%
       {arXiv:quant-ph/0011110}.

  \item[\textup{[Hil07]}]\label{ref:p15-hildebrand}
  R. Hildebrand,
  ``Positive Partial Transpose from Spectra,''
  \emph{Physical Review A} \textbf{76}, 052325 (2007).
  \href{https://doi.org/10.1103/PhysRevA.76.052325}%
       {doi:10.1103/PhysRevA.76.052325};
  \href{https://arxiv.org/abs/quant-ph/0502170}%
       {arXiv:quant-ph/0502170}.

  \item[\textup{[Joh13]}]\label{ref:p15-johnston}
  N. Johnston,
  ``Separability from Spectrum for Qubit--Qudit States,''
  \emph{Physical Review A} \textbf{88}, 062330 (2013).\newline
  \href{https://doi.org/10.1103/PhysRevA.88.062330}%
       {doi:10.1103/PhysRevA.88.062330};
  \href{https://arxiv.org/abs/1309.2006}{arXiv:1309.2006}.

  \item[\textup{[AMR+25]}]\label{ref:p15-abellanet-vidal-et-al}
  J. Abellanet-Vidal, G. Müller-Rigat, G. Rajchel-Mieldzioć, and A. Sanpera,
  ``Sufficient Criteria for Absolute Separability in Arbitrary Dimensions
  via Linear Map Inverses,''
  \emph{Reports on Progress in Physics} \textbf{88}, 107601 (2025).
  \href{https://doi.org/10.1088/1361-6633/ae0cfa}%
       {doi:10.1088/1361-6633/ae0cfa};
  \href{https://arxiv.org/abs/2410.22415}{arXiv:2410.22415}.

  \item[\textup{[AKW26]}]\label{ref:p15-ahiable-et-al}
  J. Ahiable, N. B. T. Kothakonda, and A. Winter,
  ``The Geometry of Absolute Separability and Other Convex Matrix Properties
  from Spectrum,'' arXiv:2608.03390 (2026).
  \href{https://doi.org/10.48550/arXiv.2608.03390}%
       {doi:10.48550/arXiv.2608.03390};
  \href{https://arxiv.org/abs/2608.03390}{arXiv:2608.03390}.

  \item[\textup{[Tra26]}]\label{ref:p15-tran}
  A. T. Tran,
  ``An Upper Bound for the Purity of Absolutely Positive Partial Transpose
  States,'' arXiv:2608.09832 (2026).
  \href{https://doi.org/10.48550/arXiv.2608.09832}%
       {doi:10.48550/arXiv.2608.09832};
  \href{https://arxiv.org/abs/2608.09832}{arXiv:2608.09832}.

  \item[\textup{[KW05]}]\label{ref:p15-krueger-werner}
  O. Krüger and R. F. Werner (eds.),
  ``Some Open Problems in Quantum Information Theory,''
  arXiv:quant-ph/0504166 (2005).
  \href{https://doi.org/10.48550/arXiv.quant-ph/0504166}%
       {doi:10.48550/arXiv.quant-ph/0504166};
  \href{https://arxiv.org/abs/quant-ph/0504166}%
       {arXiv:quant-ph/0504166}.
\end{enumerate}

\subsection*{Comment}

The spectral characterization was posed in the source collection
\sourcecite{ref:p15-krueger-werner}{KW05}.  The remaining gap begins when both
local dimensions are at least $3$: the complete boundary of
Eq.~\eqref{eq:p15-asep-spectra} is unknown, as is its possible equality with
Eq.~\eqref{eq:p15-appt-spectra}.

\subsection*{Tag}

Quantum separability; Partial transpose criterion; Convex geometry;
Qudit systems

\subsection*{ID}

\texttt{op\_b063bdae4363cda8}
